%% =====================================================================
%%  T-24-02  A Commitment You Did Not Choose
%%  -------------------------------------------------------------------
%%  One idea per file. Core is mandatory. Slots in canonical order.
%%  Max 700 words. No layout commands. No filler. New source material
%%  for this entry goes in sources/<BOOK>/T-24-02.tex -- never by
%%  rewriting this file (docs/RULES.md R0.1).
%% =====================================================================
%% @kind: topic
%% @id: T-24-02
%% @title: A Commitment You Did Not Choose
%% @subtitle: Recognising the position that was installed rather than decided
%% @chapter: c24
%% @order: 20
%% @status: stable
%% @sources: B4
%% @updated: 2026-09-19

\topic{T-24-02}{A Commitment You Did Not Choose}{Recognising the position that was installed rather than decided}

\begin{core}
The defence against manufactured consistency is not to resist the later request. It is to notice that the position you are defending was installed by a sequence rather than reached by a judgement, and to reopen the judgement on the current facts alone.
\end{core}
\begin{mechanism}
Consistency feels like integrity, which is why it is rarely examined. A person asked to reverse does not usually weigh the question again; they experience the request as an attack on their character and defend the position for that reason. The defence is therefore aimed at the feeling rather than at the argument.

Two moves do most of the work. The first is to distinguish reasons that existed before the commitment from reasons that appeared after it, since the second group is largely manufactured by the commitment itself. The second is to permit reversal explicitly: changing your mind on new information is a different act from being inconsistent, and the distinction has to be made consciously because the social pressure does not make it for you.
\end{mechanism}
\begin{tells}
  \item You are defending a position you cannot date --- you do not remember deciding it.
  \item Your reasons have multiplied since you agreed.
  \item Reversing feels like an admission about you rather than a calculation about the question.
  \item The original inducement has been withdrawn and the position survived it.
  \item You would not advise a friend to hold the position on these terms.
\end{tells}
\begin{moves}
  \item Ask when you took the position, and what the grounds were at that moment.
  \item List the reasons you hold it now, and mark which ones existed then.
  \item Remove the audience. Decide privately first and announce afterwards, so that identity is not part of the calculation.
  \item Reopen the question on current terms only, as though the earlier agreement had never happened.
  \item Say out loud that changing your mind on new information is not inconsistency.
\end{moves}
\begin{counters}
  \item When reversal is called weakness, agree that the position has changed and decline to defend the change; the accusation only works if you accept that consistency is what is being tested.
  \item When the original inducement is denied, do not argue about it --- note that it is no longer present and that the position was built on it.
  \item When others invoke your earlier statement, address the question rather than the record: what is being asked now, and would I agree to it today?
  \item When the cost of reversing is presented as forfeiting what you have invested, treat the investment as spent and decide on what remains.
\end{counters}
\begin{cost}
Reopening every position produces a person who cannot be relied on, and reliability is worth more than correctness in most ongoing relationships. The defence is therefore targeted rather than general: it applies where the position was installed by someone else's sequence, where the grounds have been withdrawn, or where the terms have moved. Positions genuinely reasoned through should be held, and held visibly, because that is what makes the other reversals credible.
\end{cost}

\sources{B4}
\seesources
\seealso{T-14-03, T-14-04, T-25-01, T-06-04}
